\section{Discussion}
\label{sec:discussion}

\subsection{Regime Shift vs.\ Feature Suppression}
\label{sec:mechanism}

Our linguistic analysis supports a \emph{regime shift} interpretation over a \emph{feature suppression} interpretation of iterative rejection.

{\bf What would feature suppression look like?}
If the model were learning to suppress specific detectable features through rejection feedback, we would expect AI marker words to decrease, contractions to increase, and the overall linguistic profile to converge toward human norms monotonically. None of these occur.

{\bf What do we observe instead?}
The model enters a qualitatively different generation mode. First-person pronoun usage increases 26$\times$, far exceeding human norms. The model adds conversational meta-commentary (``Alright, here's my most natural shot:'') and adopts an informal, anecdotal style. Meanwhile, AI marker words remain unchanged and contractions actually decrease. The non-monotonic trajectory---with a sharp drop at R3 and a partial rebound at R4---further suggests the model oscillates between regimes rather than converging smoothly.

This distinction matters for detector design. If evasion worked through targeted feature suppression, detectors could simply weight those features more heavily. A regime shift, however, produces a qualitatively different distribution that requires detectors to model the ``trying to sound human'' mode as a distinct class.

\subsection{Why Generic Humanization Fails}
\label{sec:generic_fails}

A surprising finding is that the naive single-step prompt (``write naturally, avoid AI patterns'') slightly \emph{increases} \binoculars detection scores (0.835 vs.\ 0.798 vanilla). We hypothesize that vague humanization instructions trigger a ``trying too hard'' mode: the model adds filler, qualifications, and an unnatural self-consciousness that is more detectable, not less.

In contrast, the strong single-step prompt succeeds because it provides \emph{specific, actionable} instructions (avoid particular words, use contractions, vary sentence length). This finding has practical implications: for users seeking to evade detectors, prompt specificity matters more than iteration count. For detector developers, this suggests that prompt-engineering-based evasion may be the more pressing threat.

\subsection{Implications}

\para{For detector trustworthiness.}
Both iterative rejection and strong single-step prompting reduce detection to near-human levels. The strong single-step approach is particularly concerning because it requires zero iteration---a single well-crafted system prompt achieves 0.342 on \binoculars, within 0.017 of the human baseline (0.325). Detector-based enforcement alone is insufficient for high-stakes decisions.

\para{For alignment research.}
The regime shift phenomenon suggests that LLMs harbor distinct generation modes activatable through conversational context. The model does not simply adjust feature weights; it switches to a different ``persona.'' This connects to work on jailbreaking and persona-based steering~\cite{madaan2023selfrefine} and suggests that conversational context exerts a qualitative, not just quantitative, influence on generation.

\para{For the \roberta detector.}
The \roberta classifier assigns human text a higher AI probability (0.320) than vanilla \gptfour output (0.264). This detector, trained on GPT-2 era text, has effectively negative discriminative power for modern LLMs. We recommend that practitioners discontinue its use for any text potentially generated by GPT-4-class models.

\subsection{Limitations}
\label{sec:limitations}

\para{Single LLM.}
We test only \gptfour. The regime shift phenomenon may differ across model families (Claude, Gemini, open-source models like LLaMA-3). Replication across models is needed.

\para{Sample size.}
Our 50-question dataset is sufficient for paired statistical tests but limits power for subgroup analyses (\eg by question type or answer length).

\para{Fixed rejection protocol.}
We use the same five escalating rejection messages for all questions. Varying the specificity of rejection feedback (\eg ``you used `delve' and `furthermore' '' vs.\ generic ``too AI'') could reveal which aspects of feedback are most effective.

\para{Detector coverage.}
We test two open-source detectors. Commercial systems (GPTZero, Turnitin) may respond differently. The \binoculars sigmoid calibration may affect absolute $P(\text{AI})$ values, though relative comparisons remain valid.

\para{Temperature and sampling.}
All generation uses temperature $= 1.0$ with seed $= 42$. Different temperature settings could change both the baseline detection rate and the effectiveness of iterative rejection.
